\documentclass[a4paper,11pt]{article}
\usepackage[utf8]{inputenc}
\usepackage[T1]{fontenc}
\usepackage[italian]{babel}
\usepackage{geometry}
\usepackage{hyperref}
\usepackage{xcolor}
\usepackage{enumitem}
\usepackage{booktabs}
\usepackage[protrusion=true,expansion=false]{microtype}

\geometry{margin=2.7cm}
\setlength{\parskip}{0.5em}
\setlength{\parindent}{0pt}

\newenvironment{citazione}{\begin{quote}\itshape}{\end{quote}}

\hypersetup{colorlinks=true, linkcolor=blue!50!black, urlcolor=blue!50!black}

\title{\textbf{The Atlas of AI}\\[0.3em]
\large Kate Crawford (2021) -- relazione di studio\\[0.3em]
\normalsize Collegamenti con il corso \emph{Etica, Società e Privacy}\\
Università di Torino -- Magistrale Informatica -- A.A. 2024/2025}
\author{}
\date{}

\begin{document}
\maketitle
\tableofcontents
\newpage

%=============================================================================
\section{Chi scrive, e con quale metodo}
%=============================================================================

Kate Crawford è ricercatrice tra Microsoft Research, la New York University e la University of Southern California, e si occupa da oltre un decennio delle implicazioni sociali, politiche e ambientali dei sistemi di intelligenza artificiale. \emph{The Atlas of AI} (Yale University Press, 2021) non è un saggio scritto da una scrivania: è il risultato di anni di viaggi e visite sul campo -- miniere di litio in Nevada, data center, archivi craniometrici, basi di lancio spaziale -- che rendono il libro insolitamente concreto rispetto alla letteratura teorica più astratta sull'etica dell'IA. È uno dei testi indicati a lezione come possibile lettura per l'esame.

Il libro si apre con un aneddoto che vale la pena ricordare per intero, perché ne anticipa la tesi. A Berlino, nei primi anni del Novecento, un cavallo chiamato Hans diventò famoso perché sembrava in grado di fare calcoli, battendo lo zoccolo il numero esatto di volte corrispondente alla soluzione di un'operazione. Un'indagine più attenta rivelò che Hans non sapeva contare: leggeva segnali corporei impercettibili e involontari del suo addestratore, che senza saperlo gli comunicava quando smettere di battere lo zoccolo. Crawford usa questa storia come metafora di un rischio che corriamo ancora oggi: scambiare per ``intelligenza'' un sistema che in realtà risponde a segnali nascosti nell'ambiente che lo circonda -- nel caso dell'IA contemporanea, enormi infrastrutture di estrazione mineraria, energia, lavoro umano e dati personali che restano invisibili a chi usa un assistente vocale o un sistema di riconoscimento facciale.

Da qui nasce la tesi che attraversa tutto il libro, riassunta in una frase che Crawford ripete quasi alla lettera: \emph{l'intelligenza artificiale non è né artificiale né intelligente}. Non è artificiale perché è fatta di materia -- minerali rari, energia, acqua, corpi di lavoratori -- ed è quindi ``fondamentalmente interconnessa alla più ampia economia estrattiva del pianeta''. Non è intelligente perché i sistemi che chiamiamo così non sono autonomi né capaci di comprendere alcunché: sono inferenza statistica addestrata su grandi quantità di dati, non pensiero. Per raccontare questa doppia smentita, Crawford scrive un libro organizzato come un atlante -- richiamando la tradizione cartografica di Aby Warburg --, cioè non una mappa neutra e universale, ma una serie di immagini parziali e situate che mettono in relazione luoghi normalmente percepiti come lontani tra loro: una miniera in Nevada, un data center in Oregon, un archivio di teschi a Philadelphia, una base di lancio in Texas. L'idea-guida che lega tutte queste immagini è che \emph{l'IA sia un registro di potere}: non crea potere dal nulla, ma registra, amplifica e rende invisibili rapporti di forza già esistenti -- economici, razziali, di genere, militari -- presentandoli come il semplice output neutro di un calcolo.

%=============================================================================
\section{Il viaggio del libro, capitolo per capitolo}
%=============================================================================

\subsection{Earth: la terra che l'IA consuma}

Il primo capitolo parte da un'osservazione tanto semplice quanto trascurata: l'intelligenza artificiale, per quanto venga venduta come qualcosa di immateriale che ``vive nel cloud'', dipende da un'industria estrattiva fatta della stessa materia di quella mineraria e petrolifera. Crawford visita Silver Peak, in Nevada, l'unica miniera di litio ancora attiva negli Stati Uniti, e mostra il contrasto stridente tra il paesaggio desertico devastato dall'estrazione e l'immagine pulita e futuristica con cui Tesla vende le proprie batterie -- arrivando a citare, quasi come una nota a margine grottesca, la proposta avanzata in passato dal fisico Edward Teller di usare ordigni nucleari per facilitare l'estrazione mineraria in quella stessa zona. Il racconto si sposta poi nella Repubblica Democratica del Congo e in Indonesia, dove l'estrazione di cobalto e stagno -- materiali indispensabili per smartphone e server -- è legata a violenza armata e a quelli che vengono chiamati ``minerali di conflitto''.

A questa geografia mineraria si aggiunge quella, altrettanto concreta, dei data center: la base della NSA a Bluffdale, nello Utah, che consuma quantità enormi d'acqua in pieno deserto, e l'impianto Google a The Dalles, in Oregon, la cui esistenza dipende da infrastrutture pubbliche sussidiate -- energia a basso costo, esenzioni fiscali -- che restano nascoste dietro la metafora immateriale del ``cloud''. Crawford cita anche uno studio sul costo energetico dell'addestramento dei grandi modelli linguistici, per mostrare quanto carbonio si nasconda dietro frasi apparentemente innocue come ``il modello ha imparato''. Il filo che lega tutti questi casi è quello che lei chiama il mito della tecnologia pulita: l'illusione che il digitale sia leggero e immateriale, quando in realtà -- per usare le sue stesse parole -- l'IA ``nasce da laghi salati in Bolivia e miniere in Congo''.

\subsection{Labor: il lavoro nascosto dietro l'automazione}

Il secondo capitolo smonta un'altra illusione, quella secondo cui i sistemi di IA funzionerebbero in modo automatico, senza bisogno di mani umane. Crawford descrive i magazzini di Amazon, dove i lavoratori sono tracciati da scanner portatili e i ritmi di lavoro sono dettati da algoritmi in una versione digitale e amplificata del taylorismo, e racconta lo sciopero organizzato nel centro logistico di Eagan, in Minnesota, da una comunità di lavoratori somalo-americani guidati da Abdi Muse e dall'organizzazione Awood Center, con uno slogan che riassume bene il punto del capitolo: ``siamo anche noi lavoratori tecnologici''.

Il caso più rilevante, però, è quello di Amazon Mechanical Turk, la piattaforma che permette di scomporre compiti complessi in microattività pagate pochi centesimi -- etichettare immagini, risolvere CAPTCHA, moderare contenuti -- svolte da milioni di persone in tutto il mondo che restano del tutto invisibili all'utente finale del servizio. Crawford riprende da Mary Gray e Sid Suri il termine \emph{ghost work} per descrivere questa forza lavoro fantasma, e da Astra Taylor il termine \emph{fauxtomation} per indicare quanto spesso ciò che viene venduto come ``automatico'' nasconda in realtà lavoro umano dietro le quinte. Anche la storia, raccontata nello stesso capitolo, della colonia industriale che Henry Ford fondò in Amazzonia -- Fordlandia -- e della rivolta operaia del 1930 contro l'imposizione dell'orologio in fabbrica, serve a ricordare che il controllo del tempo di lavoro non è un'invenzione recente, ma una continuità storica che oggi passa per il software invece che per la sirena della fabbrica. Un dettaglio linguistico che Crawford non lascia passare inosservato è che alcuni sistemi di sincronizzazione usati da Google per i propri data center, come Spanner, impiegano ancora oggi la terminologia ``master/slave'' -- un retaggio che lei legge come traccia, mai del tutto rimossa, della storia della schiavitù dentro il linguaggio tecnico dell'informatica.

\subsection{Data: dataset che non sono mai neutri}

Il terzo capitolo affronta la materia prima dell'apprendimento automatico, i dataset, mostrando che non sono mai semplice materiale grezzo, ma costruzioni storiche e politiche, spesso raccolte senza il consenso di chi vi compare. Crawford racconta la genealogia di questa pratica risalendo a Francis Galton e Alphonse Bertillon, pionieri ottocenteschi della classificazione biometrica e criminologica, diretti antenati dei moderni database di foto segnaletiche -- come quelli del NIST, costruiti a partire da detenuti che non avevano alcuna possibilità di rifiutarsi.

Il caso più importante del capitolo è quello di ImageNet, il dataset creato nel 2009 dalla ricercatrice Fei-Fei Li raschiando milioni di immagini dai motori di ricerca ed etichettandole, a un ritmo di cinquanta immagini al minuto, tramite la stessa Amazon Mechanical Turk descritta nel capitolo precedente: una dimostrazione concreta di come i due temi del libro, estrazione di lavoro ed estrazione di dati, siano in realtà la stessa cosa vista da due angolazioni diverse. ImageNet erediterà dalla struttura linguistica di WordNet, il database semantico di Princeton su cui è costruito, categorie problematiche che torneranno al centro del capitolo successivo. Crawford racconta anche una serie di casi più recenti -- dataset come DukeMTMC o MS-Celeb, costruiti raccogliendo foto pubbliche da telecamere universitarie o da celebrità senza permesso, e poi ritirati dopo scandali di privacy -- per parlare di quello che chiama ``la fine del consenso'': un'epoca in cui chiunque pubblichi una propria foto online non ha più, di fatto, alcun controllo su come quella foto potrà essere usata per addestrare un sistema di riconoscimento facciale.

\subsection{Classification: classificare è un atto di potere}

Il quarto capitolo è probabilmente il più denso di tutto il libro, e si apre con un'immagine forte: Crawford visita al Penn Museum di Philadelphia la collezione di circa mille teschi umani raccolti dal craniologo Samuel Morton, morto nel 1851, che riempiva le cavità craniche con pallini di piombo nel tentativo di ``dimostrare'' scientificamente la superiorità razziale bianca -- uno dei teschi porta semplicemente l'etichetta ``lunatico''. Crawford cita lo storico della scienza Stephen Jay Gould, che dopo aver riesaminato i dati di Morton concluse che erano un insieme di errori e manipolazioni al servizio inconsapevole di convinzioni razziste già date per scontate. È un punto di partenza scelto deliberatamente per sostenere una tesi forte: la classificazione algoritmica contemporanea non è semplicemente \emph{simile} a queste pratiche pseudoscientifiche ottocentesche, ma ne è in molti casi una continuazione diretta sotto altra forma tecnica.

Il capitolo lo dimostra ripercorrendo le categorie di ImageNet, che contava più di ventimila etichette derivate da WordNet, tra cui quasi tremila sottocategorie sotto la voce ``Persona'' -- molte delle quali, rimaste indisturbate per un decennio, erano insulti razzisti, sessisti o offensivi nei confronti delle persone con disabilità. È stato il progetto artistico \emph{ImageNet Roulette}, lanciato nel 2019 dall'artista Trevor Paglen, a rendere visibile e virale il problema, costringendo gli sviluppatori a rimuovere più della metà di quelle categorie. Crawford ricorda anche lo studio \emph{Gender Shades} di Joy Buolamwini e Timnit Gebru, che nel 2018 dimostrò come i sistemi di riconoscimento facciale di IBM, Microsoft e Amazon sbagliassero molto più spesso nel riconoscere donne con pelle scura che uomini con pelle chiara -- e la risposta poco felice di IBM, che costruì un nuovo dataset craniometrico misurando le distanze tra i punti del volto in un milione di foto, in un'eco quasi letterale del metodo di Morton.

Un altro caso citato, quello dello strumento di selezione del personale sperimentato da Amazon tra il 2014 e il 2018, è particolarmente istruttivo: il sistema continuava a penalizzare le candidate donne anche dopo che il riferimento esplicito al genere era stato rimosso dai dati, perché aveva imparato ad associare il successo professionale a verbi più comuni nei curriculum maschili. È la prova, per Crawford, che il pregiudizio non è un difetto isolabile da correggere con qualche aggiustamento tecnico: è una caratteristica intrinseca dell'atto stesso di classificare, perché ogni classificazione riflette necessariamente la storia e i pregiudizi di chi l'ha costruita.

\subsection{Affect: leggere le emozioni dai volti}

Il quinto capitolo racconta la storia, sorprendentemente fragile sul piano scientifico, dell'industria del riconoscimento delle emozioni, oggi valutata oltre diciassette miliardi di dollari. Tutto parte dal lavoro dello psicologo Paul Ekman, che nel 1967 si recò in Papua Nuova Guinea per dimostrare che le espressioni facciali delle emozioni fossero universali in tutte le culture umane. Crawford ricostruisce un dettaglio poco noto di questa storia: la ricerca di Ekman era finanziata dall'agenzia militare americana ARPA, e secondo alcuni documenti serviva anche, indirettamente, a coprire altre attività legate alle indagini del senatore Frank Church sull'uso della scienza sociale come paravento per operazioni della CIA in Cile.

Da quel lavoro nasce il Facial Action Coding System, il sistema usato ancora oggi da molte aziende per dedurre le emozioni dai movimenti del volto, applicato da società come HireVue per analizzare i colloqui video di lavoro o da Affectiva per costruire quello che definisce ``il più grande database di emozioni al mondo''. Il problema, sottolinea Crawford, è che la comunità scientifica è tutt'altro che unanime sulla possibilità di dedurre con sicurezza uno stato emotivo interiore da un'espressione facciale: una revisione del 2019 guidata dalla psicologa Lisa Feldman Barrett conclude che non è possibile inferire con sicurezza la felicità da un sorriso, la rabbia da un'espressione corrucciata, la tristezza da un broncio. L'industria, osserva Crawford con una formula efficace, ha semplicemente scelto la teoria che si adattava agli strumenti disponibili, non quella più solida. Il capitolo si chiude con un episodio politico inquietante: un blog conservatore usò diversi sistemi di riconoscimento facciale come una sorta di ``poligrafo virtuale'' per accusare la deputata Ilhan Omar di mostrare segnali di menzogna -- un esempio concreto di come queste tecnologie, prive di basi scientifiche solide, possano comunque essere usate per scopi discriminatori contro donne, e in particolare donne nere.

\subsection{State: l'intelligenza artificiale e il potere statale}

Il sesto capitolo segue il filo che collega la sorveglianza militare a quella della vita quotidiana. Crawford riprende l'archivio dei documenti diffusi da Edward Snowden per raccontare programmi della NSA come TREASUREMAP, che si proponeva di mappare l'intera rete internet, dispositivo per dispositivo, in ogni momento. Racconta poi il caso di Project Maven, il programma del Pentagono per analizzare i filmati dei drone tramite l'intelligenza artificiale di Google, che provocò una protesta interna firmata da oltre tremila dipendenti dell'azienda -- tra cui, secondo un'email diventata poi pubblica, la stessa responsabile scientifica per l'IA, Fei-Fei Li, che chiedeva di evitare a ogni costo qualsiasi menzione pubblica del coinvolgimento dell'intelligenza artificiale nel progetto. Google abbandonerà infine il contratto.

Un altro protagonista del capitolo è Palantir, l'azienda fondata da Peter Thiel che da contractor dei servizi di intelligence è diventata fornitrice di strumenti di sorveglianza per l'agenzia americana per l'immigrazione e per diversi dipartimenti di polizia locali, tra cui quello di Los Angeles. Crawford racconta anche il caso, meno conosciuto, di un punteggio sperimentale costruito da IBM durante la crisi dei rifugiati siriani del 2015 per provare a distinguere ``terroristi'' da rifugiati analizzando i dati pubblici dei social media, e quello del sistema automatico dello stato del Michigan che accusò erroneamente di frode decine di migliaia di persone che richiedevano un sussidio di disoccupazione. Il filo che lega tutti questi episodi è quello che Crawford chiama l'infiltrazione dal basso degli strumenti nati per l'intelligence militare: tecnologie pensate per i campi di battaglia o per la sorveglianza dei nemici dello stato finiscono, con sorprendente rapidità, per essere usate nelle scuole, nelle stazioni di polizia, negli uffici per la disoccupazione.

\subsection{Power: la conclusione del libro}

Nelle pagine conclusive Crawford riassume la tesi dell'intero percorso: l'intelligenza artificiale non è oggettiva, universale o neutra, ma profondamente incorporata in mondi sociali, politici, culturali ed economici specifici, costruita in modo tale da amplificare gerarchie esistenti e imporre classificazioni anguste sulla complessità della vita umana. Introduce qui un concetto che vale la pena ricordare con precisione, l'\emph{enchanted determinism}, coniato insieme allo storico Alex Campolo a partire dal caso di AlphaGo Zero, il sistema di DeepMind descritto dal suo stesso creatore come quasi ``alieno''. L'idea è che l'industria presenti i propri sistemi con un doppio registro retorico apparentemente contraddittorio: da un lato come qualcosa di incantato, magico, oltre la comprensione umana; dall'altro come qualcosa di perfettamente deterministico, capace di prevedere con assoluta certezza. Questo doppio registro, sostiene Crawford, ha l'effetto pratico di mistificare il potere reale dei sistemi di IA e di rendere quasi impossibile contestarli: se un sistema è insieme incomprensibile e infallibile, non resta spazio per il dubbio critico.

Per rendere visibile questa mistificazione, Crawford mette a confronto due immagini molto diverse: il diagramma astratto con cui viene di solito presentato AlphaGo Zero -- nessuna macchina visibile, nessun lavoro umano, nessuna impronta di carbonio -- e il progetto tecnico, molto più prosaico, del data center Google a The Dalles, che rivela tutta la sua dipendenza da infrastrutture pubbliche sussidiate. La distanza tra queste due immagini è, in fondo, il tema di tutto il libro: l'IA presentata come pura astrazione matematica nasconde sempre, da qualche parte, un'infrastruttura materiale fatta di energia, lavoro e potere. Da qui nasce anche la critica più pungente del libro nei confronti delle tante carte etiche sull'IA pubblicate da governi e aziende negli ultimi anni -- più di un centinaio solo in Europa, secondo una stima citata nel testo --, accusate di restare quasi sempre prive di qualsiasi effetto vincolante. Crawford propone di spostare l'attenzione dall'etica al potere, e rivendica quello che chiama un diritto al rifiuto: la possibilità, per una comunità, di scegliere di non usare un certo sistema di intelligenza artificiale in un certo contesto -- come hanno fatto alcune città americane bandendo il riconoscimento facciale -- citando infine la celebre frase di Audre Lorde secondo cui gli strumenti del padrone non smantelleranno mai la casa del padrone, a sottolineare che rendere l'IA più diffusa e accessibile non basta, di per sé, a renderla più giusta.

\subsection{Space: la coda del libro}

Il libro si chiude con un capitolo dedicato alla nuova corsa allo spazio dei miliardari della tecnologia, Jeff Bezos con Blue Origin ed Elon Musk con SpaceX, presentata da Crawford come l'estensione finale, e più letterale, della stessa logica estrattiva raccontata nel primo capitolo: lo spazio come ultima frontiera da colonizzare, una volta esaurite le risorse della Terra. Crawford nota con una certa amarezza che Bezos cita come proprio eroe l'ingegnere missilistico Wernher von Braun, che durante il Terzo Reich utilizzò lavoro forzato di prigionieri dei campi di concentramento per costruire i razzi V-2 -- un dettaglio storico che difficilmente compare nelle biografie celebrative della corsa allo spazio privata. Il libro termina con un episodio quasi cinematografico, raccontato in prima persona: mentre fotografa da fuori la base di lancio suborbitale di Blue Origin in Texas, Crawford viene seguita da due furgoncini neri identici -- un finale che incarna letteralmente, sulla propria pelle, il tema della sorveglianza che ha attraversato tutto il libro.

%=============================================================================
\section{Le idee che tengono insieme il libro}
%=============================================================================

Al termine del percorso capitolo per capitolo, vale la pena fissare alcune idee che ricorrono in tutto il libro e che lo rendono coerente nonostante la varietà dei casi raccontati. La prima è che terra, lavoro e dati non sono tre temi separati, ma tre facce della stessa logica estrattiva su cui si fonda l'intelligenza artificiale. La seconda è che classificare non è mai un'operazione neutra: ogni categoria imposta a un dataset porta con sé una storia, spesso una storia di potere e di esclusione, che diventa invisibile proprio perché incorporata in un'infrastruttura tecnica che usiamo ogni giorno senza pensarci. La terza è la continuità, più che la semplice somiglianza, tra le pseudoscienze ottocentesche della misurazione dei corpi -- la craniometria di Morton, l'eugenetica di Galton -- e molte pratiche di classificazione algoritmica contemporanea. La quarta è il concetto di \emph{enchanted determinism}, la doppia retorica con cui l'industria presenta l'IA come insieme magica e infallibile per sottrarla a ogni critica. La quinta è l'invisibilità del lavoro umano dietro ogni forma di apparente automazione. La sesta è la continuità tra sorveglianza militare, sorveglianza commerciale e sorveglianza amministrativa quotidiana. E infine, la settima, è il rifiuto esplicito della retorica dell'inevitabilità tecnologica: l'idea, ripetuta più volte nel libro, che il fatto che qualcosa possa essere fatto non implichi affatto che debba essere fatto.

%=============================================================================
\section{Collegamenti con il corso di Etica}
%=============================================================================

\subsection{Contro l'inevitabilismo tecnologico}

Il corso critica esplicitamente la retorica dell'inevitabilismo tecnologico, quella secondo cui l'invasività digitale sarebbe una conseguenza naturale del progresso, al pari dell'evoluzione biologica. È sorprendente quanto Crawford arrivi quasi alle stesse parole nelle pagine conclusive del libro, quando scrive che bisogna opporsi alle narrazioni dell'inevitabilità tecnologica secondo cui ``se qualcosa può essere fatto, sarà fatto'', e propone di chiedersi non dove l'IA verrà applicata, ma perché dovrebbe esserlo. Crawford va oltre, definendo l'utopismo e il distopismo tecnologico due ``gemelli metafisici'' che condividono lo stesso errore di fondo: pensare che il potere risieda nella tecnologia in sé, e non nelle strutture sociali, economiche e politiche che la producono e la usano -- un'obiezione che si lega direttamente al primo messaggio del corso, secondo cui la tecnologia e il suo impatto sulla società sono costruzioni sociali prive di inevitabilità.

\subsection{La fonte concettuale dietro il Mechanical Turk visto a lezione}

Le slide del corso usano l'immagine storica del finto automa scacchista settecentesco -- un meccanismo apparente che nascondeva un uomo al proprio interno -- per rappresentare i lavoratori della gig economy che etichettano dati restando invisibili. Il capitolo di Crawford dedicato al lavoro è, in un certo senso, la fonte concettuale più completa di quella stessa immagine: racconta nel dettaglio come funziona Amazon Mechanical Turk, introduce il concetto di \emph{ghost work} per descrivere i milioni di microlavoratori invisibili che rendono possibile l'addestramento dei sistemi di IA, e quello di \emph{fauxtomation} per indicare quanto spesso l'automazione presentata al pubblico sia in realtà lavoro umano nascosto. Lo sciopero di Eagan, in Minnesota, aggiunge a questo quadro un dettaglio che a lezione resta implicito: dietro le piattaforme che sembrano funzionare da sole ci sono comunità di lavoratori reali, capaci anche di organizzarsi e protestare.

\subsection{Un precedente storico per la critica al bias dei modelli generativi}

A lezione viene mostrato come Stable Diffusion, alla richiesta di generare l'immagine di un giudice, produca quasi solo uomini bianchi, traendone la conclusione che il problema non sia tecnico ma politico, perché le scelte alla base di quel risultato le fanno gli ingegneri di un'azienda. Il progetto \emph{ImageNet Roulette} dell'artista Trevor Paglen, raccontato nel capitolo di Crawford sulla classificazione, è il precedente storico quasi esatto di questo tipo di critica: un'opera artistica che ha reso visibile al pubblico, prima ancora che ai tecnici, quanto le categorie di un dataset di addestramento potessero contenere pregiudizi razzisti e sessisti rimasti indisturbati per anni. Il caso dello strumento di selezione del personale di Amazon, che continuava a discriminare le candidate donne anche dopo la rimozione del riferimento esplicito al genere, rafforza ulteriormente questo punto: il pregiudizio non è un bug isolabile, ma una caratteristica strutturale di ogni atto di classificazione.

\subsection{Classificare come esercizio di potere, non come fatto neutro}

Il costruttivismo sociale di Searle, discusso a lezione, insegna che la realtà sociale in cui viviamo è fatta di regole costitutive condivise, non di fatti bruti della natura. Il capitolo di Crawford sulla classificazione applica una logica complementare al potere di classificare: una categoria di ImageNet, una volta incorporata nell'infrastruttura tecnica che usiamo ogni giorno, diventa quasi invisibile senza perdere nulla del proprio potere -- esattamente il modo in cui i fatti istituzionali di Searle smettono di essere percepiti come costruzioni, pur continuando a strutturare la vita quotidiana di chi vi è sottoposto.

\subsection{Dalla sorveglianza militare a quella di tutti i giorni}

Il capitolo di Crawford sullo Stato offre molto materiale utile per arricchire il discorso del corso sulla sorveglianza di massa: dai programmi della NSA, a Palantir e il suo lavoro per l'agenzia americana per l'immigrazione, fino al caso del riconoscimento delle emozioni usato contro la deputata Ilhan Omar. L'idea che strumenti nati per l'intelligence militare finiscano per filtrare verso il basso, fino alle scuole e agli uffici per la disoccupazione, aggiunge una dimensione storica e genealogica a un tema che a lezione viene introdotto soprattutto attraverso esempi più recenti, come i sistemi di sorveglianza di massa delle minoranze.

\subsection{Una lista degli appunti che è quasi un riassunto del libro}

Tra gli appunti del corso compare un elenco intitolato ``cosa non viene mostrato'' dell'intelligenza artificiale: l'inquinamento delle miniere di litio, le proteste delle popolazioni la cui vita viene distrutta dall'estrazione mineraria, le terre rare estratte in paesi poveri e in zone di guerra spesso da manodopera minorile, le condizioni quasi schiavili in cui centinaia di migliaia di persone lavorano alla produzione dei dispositivi elettronici, gli enormi data center alimentati da fonti non rinnovabili, le foto usate per addestrare i sistemi senza che chi le ha scattate o chi vi compare ne abbia dato il permesso, e i lavoratori sottopagati che, come nel documentario \emph{The Cleaners}, sono costretti a visionare i contenuti più violenti per ``pulire'' social media e dataset. È difficile non notare che questa lista corrisponde quasi punto per punto ai primi due capitoli di \emph{The Atlas of AI}: Silver Peak e le miniere di litio, il cobalto del Congo e le terre rare dell'Indonesia raccontate in \emph{Earth}, e Amazon Mechanical Turk, il \emph{ghost work} e l'ossessione per la produttività erede del taylorismo raccontati in \emph{Labor}. Il libro di Crawford fornisce, in un certo senso, lo sviluppo narrativo e documentato in profondità di ciò che a lezione viene presentato come un elenco sintetico.

\subsection{Tecnofeudalesimo e la nuova corsa allo spazio}

Il corso introduce il concetto di \emph{tecnofeudalesimo} a partire da Amazon: Jeff Bezos non si comporta più come un imprenditore che compete in un mercato, ma come un aristocratico che vive di rendita sul territorio digitale che possiede, imponendo le proprie regole a chi lo attraversa. Tra gli esempi citati a lezione di come questa logica di erosione dei beni comuni si estenda anche al di fuori della Terra compaiono i satelliti Starlink in orbita bassa. È esattamente il punto in cui arriva Crawford nella coda del libro, dedicata alla nuova corsa allo spazio di Bezos e Musk: lo spazio descritto come un'ultima frontiera da colonizzare una volta esaurite le risorse terrestri, con la stessa logica estrattiva e la stessa concentrazione di potere privato raccontate nel resto del libro -- al punto che Crawford nota con amarezza come Bezos citi come proprio eroe l'ingegnere Wernher von Braun, che usò lavoro forzato per costruire i razzi nazisti.

\subsection{Cartesio, La Mettrie e il corpo che l'intelligenza artificiale vorrebbe far dimenticare}

A lezione viene ricostruita la svolta di Cartesio, che rese la realtà materiale -- compreso il corpo umano -- oggetto di scienza, ma lasciò la mente in una posizione separata e quasi divina, il celebre dualismo cartesiano; e poi la radicalizzazione di questa idea nell'\emph{Homme Machine} di Julien Offray de La Mettrie, secondo cui anche i processi mentali sarebbero spiegabili con principi fisici e meccanici, senza alcun bisogno di un'anima immateriale. È una cornice filosofica che illumina bene il punto di partenza di Crawford: l'affermazione che l'intelligenza artificiale non sia né artificiale né intelligente è, in fondo, un rovesciamento del dualismo cartesiano applicato alla macchina invece che all'uomo. L'industria dell'IA preferisce presentare i propri sistemi come pura mente che calcola, quasi immateriale -- è il concetto di \emph{enchanted determinism} descritto nella conclusione del libro -- mentre Crawford insiste nel ricordare che dietro quella mente apparente ci sono sempre un corpo, una miniera, un magazzino: la materialità che il dualismo cartesiano, e oggi la retorica dell'IA, preferirebbero lasciare in ombra.

\subsection{Dati come bene comune eroso senza consenso}

Il corso distingue un \emph{social network}, che serve a connettere le persone, da un \emph{social media}, che serve a distribuire contenuti, osservando come molti social network siano progressivamente diventati social media -- un'evoluzione riassunta dalla frase ``se non stai pagando per il prodotto, sei tu il prodotto''. È la stessa dinamica di estrazione, applicata ai contenuti invece che ai dati di addestramento, che Crawford descrive nel capitolo sui dati a proposito della cosiddetta ``fine del consenso'': foto pubblicate online per restare in contatto con amici e parenti che finiscono, senza che nessuno lo abbia davvero deciso, dentro dataset usati per addestrare sistemi di riconoscimento facciale. In entrambi i casi una risorsa che nasce come bene condiviso tra persone -- relazioni sociali, fotografie di famiglia -- viene silenziosamente trasformata in materia prima per un'attività commerciale di cui chi l'ha prodotta non è mai stato informato.

%=============================================================================
\section{Tabella di sintesi}
%=============================================================================

\begin{table}[h!]
\centering
\renewcommand{\arraystretch}{1.35}
\begin{tabular}{p{6cm}p{8.5cm}}
\toprule
\textbf{Concetto del corso di Etica} & \textbf{Come si ritrova in \emph{The Atlas of AI}} \\
\midrule
Inevitabilismo tecnologico come retorica & La ``politics of refusal'' contro la narrazione ``se può essere fatto, sarà fatto'' \\
Costruzione sociale della tecnologia & L'IA come ``registro di potere'', non come output neutro di un calcolo \\
Costruttivismo sociale di Searle & Le categorie di un dataset come fatti istituzionali, non fatti bruti, rese invisibili dall'infrastruttura tecnica \\
Bias dell'AI come scelta politica & ImageNet Roulette e il caso Amazon: il pregiudizio è una caratteristica della classificazione, non un bug \\
Sfruttamento del lavoro invisibile & Amazon Mechanical Turk, ghost work, fauxtomation, lo sciopero di Eagan \\
Sorveglianza di massa & NSA, Project Maven, Palantir, il caso Ilhan Omar \\
Lista ``cosa non viene mostrato'' dell'AI & Corrisponde quasi punto per punto ai capitoli Earth (miniere, terre rare) e Labor (Mechanical Turk, The Cleaners) \\
Tecnofeudalesimo e Starlink & La coda \emph{Space}: la corsa allo spazio di Bezos e Musk come estensione della logica estrattiva \\
Dualismo cartesiano e L'Homme Machine & ``Enchanted determinism'': l'industria nasconde il corpo materiale dell'IA dietro l'illusione di una mente pura \\
Social media come estrazione senza consenso & Il capitolo Data e ``la fine del consenso'': foto pubbliche diventate materiale di addestramento senza permesso \\
\bottomrule
\end{tabular}
\caption{Sintesi dei collegamenti tra il corso di Etica e \emph{The Atlas of AI}.}
\end{table}

%=============================================================================
\section{Domande utili per l'esposizione orale}
%=============================================================================

Una traccia di domande con cui mettersi alla prova, utile per ripassare prima dell'esame:

\begin{enumerate}
  \item Spiega il senso della frase ``l'intelligenza artificiale non è né artificiale né intelligente''. In che modo questa tesi struttura i capitoli su terra, lavoro e dati?
  \item Cos'è l'\emph{enchanted determinism}? Usa il caso di AlphaGo Zero per spiegare perché questa doppia retorica rende difficile criticare i sistemi di intelligenza artificiale.
  \item In che senso Crawford parla di continuità, e non solo di analogia, tra la craniometria di Samuel Morton e i moderni dataset di addestramento come ImageNet?
  \item Spiega i concetti di \emph{ghost work} e \emph{fauxtomation} a partire dal caso di Amazon Mechanical Turk e dello sciopero di Eagan, in Minnesota.
  \item Perché Crawford sostiene, a partire dal caso dello strumento di selezione del personale di Amazon, che il pregiudizio algoritmico sia una caratteristica della classificazione stessa e non un semplice errore tecnico correggibile?
  \item Come si collega il capitolo sullo Stato alla tesi più generale del libro secondo cui l'IA è un ``registro di potere''?
  \item Cosa intende Crawford quando propone di spostare l'attenzione dall'etica dell'IA al potere dell'IA? Perché ritiene insufficienti le carte etiche pubblicate da aziende e governi?
  \item Collega il dualismo cartesiano e \emph{L'Homme Machine} di La Mettrie, visti a lezione, al concetto di \emph{enchanted determinism}: in che senso l'industria dell'IA preferisce nascondere la materialità del corpo che la sostiene?
  \item In che modo la coda del libro dedicata alla corsa allo spazio di Bezos e Musk si collega al concetto di tecnofeudalesimo discusso a lezione?
\end{enumerate}

Un'ultima nota utile per l'orale: \emph{The Atlas of AI} è particolarmente efficace per arricchire con dettagli concreti -- nomi, date, citazioni testuali -- argomenti già discussi a lezione, in particolare il riferimento al Mechanical Turk e al bias dei modelli generativi, che nel libro trovano un trattamento storico molto più ampio. È un buon modo per mostrare approfondimento autonomo durante l'esposizione orale.

\end{document}
